\documentclass{amsart}
\usepackage{amsmath,amssymb,amsthm,hyperref}
\usepackage{algorithm}
\usepackage{algpseudocode}

\newtheorem{theorem}{Theorem}
\newtheorem{lemma}{Lemma}
\newtheorem{proposition}{Proposition}
\newtheorem{corollary}{Corollary}
\theoremstyle{definition}
\newtheorem{definition}{Definition}
\newtheorem{remark}{Remark}

\newcommand{\Chat}{\widehat{C}}
\newcommand{\Lhat}{\widehat{L}}
\newcommand{\Nhat}{\widehat{N}}
\newcommand{\E}{\mathbb{E}}
\renewcommand{\Pr}{\mathbb{P}}

\begin{document}

\title{Measurement, estimation, and provable mapping quality in heterogeneous computing}
\maketitle

\section{Overview and reduction to a clean mathematical core}

The problem asks us to (a) define computable runtime descriptors $L_i(\tau)$ and
$N_{ij}(\tau)$, (b) produce completion-time estimators $\Chat$ with controlled
accuracy, and (c) prove that minimizing the \emph{estimated} objective yields a
mapping provably close to the optimum of the \emph{exact} objective.

The engineering content of (a) and (b) (which descriptors to read, how to fit a
predictor) is real but is not where the mathematical difficulty lies. The
rigorous, problem-defining content is the following implication, which we isolate
and prove in full generality:

\begin{quote}
\emph{If the estimator is accurate in a multiplicative (relative-error) sense
uniformly over all task/subtask--machine pairs, then any mapping that exactly
minimizes the estimated objective is, under the exact objective, within a factor
$\frac{1+\varepsilon}{1-\varepsilon}$ of optimal, for a broad and explicitly
characterized class of objectives (makespan and weighted average completion time
among them).}
\end{quote}

We then close the loop by exhibiting concrete measurement and estimation
procedures (averaging of runtime samples) for which the required multiplicative
accuracy holds with high probability, with an explicit sample-complexity bound.
Sections~\ref{sec:model}--\ref{sec:transfer} give the mathematics; the
construction is stated explicitly in Algorithms~\ref{alg:descr} and
\ref{alg:map}, with proofs of correctness, accuracy, and complexity.

Throughout, ``mapping'' means a function assigning each schedulable unit (task or
subtask) to a machine; the completion-time function already incorporates load,
communication, and precedence through its dependence on $(L,N)$, so we treat it
as a primitive once $(L,N)$ is fixed, exactly as in the problem statement.

\section{Model and definitions}\label{sec:model}

\begin{definition}[Schedulable units and mappings]\label{def:units}
Let $U$ be the finite set of \emph{schedulable units}: $U = T \cup S$, where
$T=\{t_1,\dots,t_k\}$ are tasks and $S$ is the (finite) set of all subtasks
obtained by decomposing tasks. Let $M=\{m_1,\dots,m_n\}$ be the machines. A
\emph{mapping} is a function $x\colon U \to M$. We write $x_u=i$ if unit $u$ is
assigned to machine $m_i$, and we write $\mathcal{X}$ for the (finite) set of all
mappings consistent with the problem's hard constraints (e.g.\ co-location or
exclusivity constraints); $\mathcal X$ is required only to be nonempty.
\end{definition}

\begin{definition}[Completion-time matrix]\label{def:cmatrix}
Fix a load/status profile $(L,N)$, where $L=(L_1,\dots,L_n)$ collects the machine
descriptors and $N=(N_{ij})$ the link descriptors. The exact completion-time
function $C(\cdot,\cdot;L,N)$ from the problem induces a matrix
$C \in \mathbb{R}_{>0}^{U\times M}$ with entries $C_{ui} := C(u,m_i;L,N)$,
the predicted completion time of unit $u$ on machine $m_i$ under profile $(L,N)$.
We assume $C_{ui}>0$ for all admissible $(u,i)$, which holds because any
nonempty unit takes strictly positive time. An \emph{estimator} is a matrix
$\Chat\in\mathbb{R}_{>0}^{U\times M}$ with entries
$\Chat_{ui}=\Chat(u,m_i;L,N)$.
\end{definition}

\begin{definition}[Objective]\label{def:obj}
An \emph{objective} is a function
$F\colon \mathbb{R}_{>0}^{U\times M}\times\mathcal{X}\to\mathbb{R}_{\ge 0}$,
written $F(D,x)$, that maps a completion-time matrix $D$ and a mapping $x$ to a
cost. The two canonical objectives are:
\begin{align}
\text{(makespan)}\quad
F_{\max}(D,x) &:= \max_{i\in[n]} \sum_{u:\,x_u=i} D_{ui}, \label{eq:makespan}\\
\text{(weighted avg.\ completion)}\quad
F_{w}(D,x) &:= \sum_{u\in U} w_u\, D_{u,x_u}, \label{eq:wsum}
\end{align}
where $w_u\ge 0$ are fixed weights.
\end{definition}

We now extract the exact structural property of $F$ that makes the transfer
theorem work.

\begin{definition}[Scale-coherence]\label{def:coherent}
An objective $F$ is \emph{scale-coherent} if for every constant $c>0$, every
matrix $D\in\mathbb{R}_{>0}^{U\times M}$, and every mapping $x\in\mathcal{X}$,
\begin{equation}\label{eq:scale}
F(cD, x) = c\, F(D,x),
\end{equation}
and if $F$ is \emph{monotone}, i.e.\ $D\le D'$ entrywise implies
$F(D,x)\le F(D',x)$ for every $x$.
\end{definition}

\begin{lemma}[The canonical objectives are scale-coherent]\label{lem:canonical}
Both $F_{\max}$ in \eqref{eq:makespan} and $F_w$ in \eqref{eq:wsum} are
scale-coherent.
\end{lemma}

\begin{proof}
\emph{Homogeneity \eqref{eq:scale}.} For $F_{\max}$ and any $c>0$,
\[
F_{\max}(cD,x)=\max_{i}\sum_{u:x_u=i} c\,D_{ui}
=c\max_{i}\sum_{u:x_u=i} D_{ui}=c\,F_{\max}(D,x),
\]
using that $c>0$ may be pulled out of a finite sum and out of a maximum of
nonnegative numbers. For $F_w$,
$F_w(cD,x)=\sum_u w_u\, c D_{u,x_u}=c\sum_u w_u D_{u,x_u}=c\,F_w(D,x)$.

\emph{Monotonicity.} If $D\le D'$ entrywise then for fixed $x$ and each $i$,
$\sum_{u:x_u=i}D_{ui}\le\sum_{u:x_u=i}D'_{ui}$ (sum of termwise inequalities of
nonnegative numbers), and the maximum over $i$ of larger quantities is larger, so
$F_{\max}(D,x)\le F_{\max}(D',x)$. For $F_w$, since $w_u\ge0$,
$w_u D_{u,x_u}\le w_u D'_{u,x_u}$ termwise, hence $F_w(D,x)\le F_w(D',x)$.
\end{proof}

\begin{remark}
Scale-coherence is a genuinely mild condition: it holds for any objective that is
a monotone, positively homogeneous-of-degree-one aggregate of the per-unit
completion times. This includes $\ell_p$ load norms
$\big(\sum_i (\sum_{u:x_u=i} D_{ui})^p\big)^{1/p}$ ($p\ge1$), maximum weighted
completion time $\max_u w_u D_{u,x_u}$, and convex combinations of
scale-coherent objectives. The proof below uses \emph{only}
Definition~\ref{def:coherent}, so it applies verbatim to all of these.
\end{remark}

\section{Measurement: computing the descriptors (part (a))}\label{sec:measure}

We make the descriptors concrete and computable from runtime samples.

\begin{definition}[Descriptors]\label{def:descriptors}
At sampling instant $\tau$ each machine $m_i$ exports a fixed-dimensional vector
\[
L_i(\tau)=\big(q_i(\tau),\, \rho^{\mathrm{cpu}}_i(\tau),\, \rho^{\mathrm{mem}}_i(\tau),\,
a_i(\tau),\, \theta_i\big)\in\mathbb{R}^{5},
\]
where $q_i$ is queue length (number of enqueued units), $\rho^{\mathrm{cpu}}_i,
\rho^{\mathrm{mem}}_i\in[0,1]$ are utilizations, $a_i\in\{0,1\}$ is availability,
and $\theta_i>0$ is the (static) nominal throughput of $m_i$. Each link exports
\[
N_{ij}(\tau)=\big(\beta_{ij}(\tau),\,\lambda_{ij}(\tau)\big)\in\mathbb{R}^2_{\ge0},
\]
effective bandwidth and latency.
\end{definition}

These raw readings are noisy. We define the \emph{computed} descriptor as the
empirical mean over a measurement window, which is what Algorithm~\ref{alg:descr}
returns and what the estimators consume.

\begin{algorithm}[h]
\caption{ComputeDescriptors}\label{alg:descr}
\begin{algorithmic}[1]
\Require window length $r\in\mathbb{N}$; raw sample streams $L_i(\tau_1),\dots,L_i(\tau_r)$ and $N_{ij}(\tau_1),\dots,N_{ij}(\tau_r)$ for all $i,j$.
\Ensure averaged descriptors $\bar L_i,\ \bar N_{ij}$.
\For{$i=1$ \textbf{to} $n$}
  \State $\bar L_i \gets \frac{1}{r}\sum_{p=1}^{r} L_i(\tau_p)$
  \For{$j=1$ \textbf{to} $n$, $j\ne i$}
     \State $\bar N_{ij} \gets \frac{1}{r}\sum_{p=1}^{r} N_{ij}(\tau_p)$
  \EndFor
\EndFor
\State \Return $(\bar L_i)_i,\ (\bar N_{ij})_{ij}$
\end{algorithmic}
\end{algorithm}

\noindent\textbf{Correctness, termination, complexity.} The algorithm performs a
fixed number of additions and divisions: termination is immediate. It computes
exactly the empirical means, so it is correct by definition. It runs in
$O(r\,n^2)$ arithmetic operations and $O(n^2)$ space (storing the running sums).

\section{Estimation accuracy from measurement (part (b))}\label{sec:estimate}

We now give a procedure whose output $\Chat$ satisfies the multiplicative
accuracy that Section~\ref{sec:transfer} requires, and we bound its sample
complexity. The estimator is the standard ``measure repeatedly and average''
predictor, instantiated per (unit, machine) pair.

\begin{definition}[Sampled completion times]\label{def:samples}
Fix a unit $u$ and machine $m_i$ under a stationary profile $(L,N)$. Let
$Y^{(1)}_{ui},\dots,Y^{(r)}_{ui}$ be $r$ independent measured completion times of
$u$ on $m_i$ obtained from runtime probes or a calibrated performance model fed
with the descriptors of Section~\ref{sec:measure}. Assume:
\begin{enumerate}
\item[(M1)] (Unbiasedness) $\E[Y^{(p)}_{ui}]=C_{ui}$ for all $p$;
\item[(M2)] (Boundedness) $0\le Y^{(p)}_{ui}\le b_{ui}$ almost surely, for a known
$b_{ui}>0$;
\item[(M3)] (Positivity) $C_{ui}\ge a_{ui}>0$ for a known lower bound $a_{ui}$.
\end{enumerate}
Define the estimator $\Chat_{ui} := \frac1r\sum_{p=1}^{r} Y^{(p)}_{ui}$.
\end{definition}

Assumptions (M1)--(M3) are the precise, checkable hypotheses under which a
relative-error guarantee is attainable: (M1) says probing is calibrated, (M2)
says completion times are bounded (true on real hardware: a unit cannot exceed a
timeout $b_{ui}$), and (M3) says nonempty units take bounded-away-from-zero time.

\begin{proposition}[Sample complexity for multiplicative accuracy]\label{prop:samples}
Fix $\varepsilon\in(0,1)$ and $\delta\in(0,1)$. Under (M1)--(M3), if the number of
samples satisfies
\begin{equation}\label{eq:rbound}
r \ \ge\ \frac{b_{ui}^2}{2\,a_{ui}^2\,\varepsilon^2}\,\ln\!\frac{2}{\delta},
\end{equation}
then $\Pr\!\big[\,(1-\varepsilon)\,C_{ui}\le \Chat_{ui}\le (1+\varepsilon)\,C_{ui}\,\big]\ge 1-\delta.$
\end{proposition}

\begin{proof}
By (M1), $\E[\Chat_{ui}]=\frac1r\sum_p \E[Y^{(p)}_{ui}]=C_{ui}$. The summands are
independent and, by (M2), lie in $[0,b_{ui}]$. Hoeffding's inequality
(for independent $X_p\in[0,b]$, $\Pr[|\frac1r\sum X_p-\mu|\ge s]\le
2\exp(-2rs^2/b^2)$) with $s=\varepsilon a_{ui}$ gives
\[
\Pr\!\big[\,|\Chat_{ui}-C_{ui}|\ge \varepsilon a_{ui}\,\big]
\ \le\ 2\exp\!\Big(-\frac{2 r \varepsilon^2 a_{ui}^2}{b_{ui}^2}\Big).
\]
The right side is $\le \delta$ exactly when $r\ge \frac{b_{ui}^2}{2a_{ui}^2\varepsilon^2}\ln\frac2\delta$,
which is \eqref{eq:rbound}. On the complementary event,
$|\Chat_{ui}-C_{ui}|<\varepsilon a_{ui}\le \varepsilon C_{ui}$ by (M3), i.e.\
$(1-\varepsilon)C_{ui}<\Chat_{ui}<(1+\varepsilon)C_{ui}$.
\end{proof}

\begin{corollary}[Uniform multiplicative accuracy]\label{cor:uniform}
Let $P$ be the number of admissible $(u,i)$ pairs. If for each pair the estimator
uses $r_{ui}\ge \frac{b_{ui}^2}{2a_{ui}^2\varepsilon^2}\ln\frac{2P}{\delta}$
independent samples, then with probability at least $1-\delta$,
\begin{equation}\label{eq:Eeps}
(1-\varepsilon)\,C_{ui}\ \le\ \Chat_{ui}\ \le\ (1+\varepsilon)\,C_{ui}
\qquad\text{for all admissible }(u,i).
\end{equation}
\end{corollary}

\begin{proof}
Apply Proposition~\ref{prop:samples} with $\delta$ replaced by $\delta/P$ to each
of the $P$ pairs; each fails with probability at most $\delta/P$. By the union
bound the probability that any pair fails is at most $P\cdot(\delta/P)=\delta$, so
\eqref{eq:Eeps} holds for all pairs simultaneously with probability $\ge1-\delta$.
\end{proof}

We call event \eqref{eq:Eeps} the \emph{$\varepsilon$-accuracy event}
$E_\varepsilon$. Corollary~\ref{cor:uniform} shows $E_\varepsilon$ is attainable
with explicitly bounded measurement effort. The expected relative error of each
entry is also controlled: $\E\big[|\Chat_{ui}-C_{ui}|\big]/C_{ui}
\le \sqrt{\operatorname{Var}(\Chat_{ui})}/C_{ui}
\le b_{ui}/(2 a_{ui}\sqrt{r_{ui}})$ by (M2)--(M3) and the variance bound for
bounded variables, so it is $O(1/\sqrt{r})$.

\section{Transfer theorem: accuracy $\Rightarrow$ mapping quality (part (c))}\label{sec:transfer}

This is the core result. We work \emph{deterministically} on the event
$E_\varepsilon$ of \eqref{eq:Eeps}; combining with Corollary~\ref{cor:uniform}
then yields a high-probability statement.

\begin{lemma}[Pointwise objective sandwich]\label{lem:sandwich}
Let $F$ be scale-coherent (Definition~\ref{def:coherent}). Suppose $\Chat$ and
$C$ satisfy \eqref{eq:Eeps} with $\varepsilon\in(0,1)$. Then for every mapping
$x\in\mathcal{X}$,
\begin{equation}\label{eq:sandwich}
(1-\varepsilon)\,F(C,x)\ \le\ F(\Chat,x)\ \le\ (1+\varepsilon)\,F(C,x).
\end{equation}
\end{lemma}

\begin{proof}
By \eqref{eq:Eeps}, $\Chat \le (1+\varepsilon)C$ entrywise. By monotonicity and
then homogeneity \eqref{eq:scale} of $F$,
\[
F(\Chat,x)\ \le\ F\big((1+\varepsilon)C,\,x\big)\ =\ (1+\varepsilon)\,F(C,x),
\]
which is the right inequality. Likewise $\Chat\ge(1-\varepsilon)C$ entrywise, and
since $1-\varepsilon>0$ the matrix $(1-\varepsilon)C$ is admissible
(positive entries), so monotonicity and homogeneity give
\[
F(\Chat,x)\ \ge\ F\big((1-\varepsilon)C,\,x\big)\ =\ (1-\varepsilon)\,F(C,x),
\]
the left inequality.
\end{proof}

\begin{theorem}[Approximation guarantee under estimated optimization]\label{thm:main}
Let $F$ be scale-coherent. Suppose the $\varepsilon$-accuracy event
$E_\varepsilon$ of \eqref{eq:Eeps} holds with $\varepsilon\in(0,1)$. Let
\[
\hat{x}\ \in\ \arg\min_{x\in\mathcal{X}} F(\Chat,x)
\qquad\text{and}\qquad
x^\star\ \in\ \arg\min_{x\in\mathcal{X}} F(C,x).
\]
Then $\hat{x}$ is a near-optimal mapping for the \emph{exact} objective:
\begin{equation}\label{eq:approx}
F(C,\hat{x})\ \le\ \frac{1+\varepsilon}{1-\varepsilon}\;F(C,x^\star).
\end{equation}
\end{theorem}

\begin{proof}
Both $\arg\min$ sets are nonempty because $\mathcal X$ is finite and nonempty.
Apply the left inequality of \eqref{eq:sandwich} to $\hat x$ and the right
inequality to $x^\star$:
\[
(1-\varepsilon)\,F(C,\hat x)\ \le\ F(\Chat,\hat x),
\qquad
F(\Chat,x^\star)\ \le\ (1+\varepsilon)\,F(C,x^\star).
\]
By optimality of $\hat x$ for $F(\Chat,\cdot)$, and since $x^\star\in\mathcal X$,
\[
F(\Chat,\hat x)\ \le\ F(\Chat,x^\star).
\]
Chaining the three displays,
\[
(1-\varepsilon)\,F(C,\hat x)\ \le\ F(\Chat,\hat x)\ \le\ F(\Chat,x^\star)
\ \le\ (1+\varepsilon)\,F(C,x^\star).
\]
Dividing by $1-\varepsilon>0$ gives \eqref{eq:approx}.
\end{proof}

\begin{corollary}[End-to-end high-probability guarantee]\label{cor:e2e}
Fix $\varepsilon\in(0,1)$, $\delta\in(0,1)$, and a scale-coherent objective $F$.
Run \textsc{ComputeDescriptors} (Algorithm~\ref{alg:descr}), build the estimator
$\Chat$ as in Definition~\ref{def:samples} with
$r_{ui}\ge \frac{b_{ui}^2}{2a_{ui}^2\varepsilon^2}\ln\frac{2P}{\delta}$ samples per
pair, and output any
$\hat x\in\arg\min_{x\in\mathcal X}F(\Chat,x)$ (Algorithm~\ref{alg:map}). Then,
with probability at least $1-\delta$,
\[
F(C,\hat x)\ \le\ \frac{1+\varepsilon}{1-\varepsilon}\,\min_{x\in\mathcal X}F(C,x).
\]
\end{corollary}

\begin{proof}
By Corollary~\ref{cor:uniform}, $E_\varepsilon$ holds with probability
$\ge1-\delta$. On $E_\varepsilon$, Theorem~\ref{thm:main} gives the stated bound.
\end{proof}

\begin{theorem}[Robustness to approximate optimization]\label{thm:approxalg}
Let $F$ be scale-coherent and suppose $E_\varepsilon$ holds with
$\varepsilon\in(0,1)$. Let $\gamma\ge 1$ and let $\tilde x\in\mathcal X$ be any
$\gamma$-approximate minimizer of the estimated objective, i.e.
\[
F(\Chat,\tilde x)\ \le\ \gamma\,\min_{x\in\mathcal X} F(\Chat,x).
\]
Then $\tilde x$ satisfies, under the exact objective,
\begin{equation}\label{eq:approxalg}
F(C,\tilde x)\ \le\ \gamma\,\frac{1+\varepsilon}{1-\varepsilon}\;\min_{x\in\mathcal X}F(C,x).
\end{equation}
\end{theorem}

\begin{proof}
Let $x^\star\in\arg\min_x F(C,x)$ (nonempty since $\mathcal X$ is finite,
nonempty). By the left inequality of \eqref{eq:sandwich} applied to $\tilde x$,
the $\gamma$-approximation hypothesis, the trivial bound
$\min_x F(\Chat,x)\le F(\Chat,x^\star)$, and the right inequality of
\eqref{eq:sandwich} applied to $x^\star$,
\[
(1-\varepsilon)F(C,\tilde x)\le F(\Chat,\tilde x)
\le \gamma\min_x F(\Chat,x)
\le \gamma F(\Chat,x^\star)
\le \gamma(1+\varepsilon)F(C,x^\star).
\]
Dividing by $1-\varepsilon>0$ yields \eqref{eq:approxalg}. Taking $\gamma=1$
recovers Theorem~\ref{thm:main}.
\end{proof}

\begin{remark}[Tightness of the factor $\tfrac{1+\varepsilon}{1-\varepsilon}$]\label{rem:tight}
The bound in Theorem~\ref{thm:main} cannot be improved for the class of
scale-coherent objectives. Take $U=\{u\}$, $M=\{m_1,m_2\}$, $F=F_{\max}$, exact
times $C_{u1}=1-\varepsilon$, $C_{u2}=1+\varepsilon$, and estimates
$\Chat_{u1}=(1+\varepsilon)C_{u1}=(1+\varepsilon)(1-\varepsilon)=\Chat_{u2}
=(1-\varepsilon)C_{u2}$. Both estimates are equal, so $\hat x$ may assign $u$ to
$m_1$, giving $F(C,\hat x)=1+\varepsilon$, whereas the true optimum is
$F(C,x^\star)=1-\varepsilon$; the ratio is exactly
$\frac{1+\varepsilon}{1-\varepsilon}$. Both estimates obey \eqref{eq:Eeps}, so the
constant is attained.
\end{remark}

\begin{algorithm}[h]
\caption{EstimatedOptimalMapping}\label{alg:map}
\begin{algorithmic}[1]
\Require accuracy $\varepsilon\in(0,1)$, confidence $\delta\in(0,1)$, scale-coherent objective $F$, sample bounds $a_{ui},b_{ui}$, mapping family $\mathcal X$.
\Ensure mapping $\hat x$ with $F(C,\hat x)\le \frac{1+\varepsilon}{1-\varepsilon}\min_{x}F(C,x)$ w.p.\ $\ge 1-\delta$.
\State $P\gets$ number of admissible $(u,i)$ pairs
\For{each admissible pair $(u,i)$}
  \State $r_{ui}\gets \big\lceil \frac{b_{ui}^2}{2a_{ui}^2\varepsilon^2}\ln\frac{2P}{\delta}\big\rceil$
  \State draw samples $Y^{(1)}_{ui},\dots,Y^{(r_{ui})}_{ui}$ satisfying (M1)--(M3)
  \State $\Chat_{ui}\gets \frac{1}{r_{ui}}\sum_{p=1}^{r_{ui}} Y^{(p)}_{ui}$
\EndFor
\State $\hat x \gets \arg\min_{x\in\mathcal X} F(\Chat,x)$
\State \Return $\hat x$
\end{algorithmic}
\end{algorithm}

\noindent\textbf{Correctness.} Lines 2--6 realize the estimator of
Definition~\ref{def:samples} at the sample size of
Corollary~\ref{cor:uniform}, so $E_\varepsilon$ holds w.p.\ $\ge1-\delta$; line~7
returns an exact minimizer of $F(\Chat,\cdot)$. Corollary~\ref{cor:e2e} then
certifies the output guarantee. \textbf{Termination:} all loops are finite and
$\mathcal X$ is finite. \textbf{Complexity:} sampling costs
$\sum_{u,i} r_{ui} = O\!\big(P\,\frac{b_{\max}^2}{a_{\min}^2\varepsilon^2}\ln\frac{P}{\delta}\big)$
probes (here $b_{\max}=\max b_{ui}$, $a_{\min}=\min a_{ui}$); the final
optimization in line~7 costs the time to minimize $F$ over $\mathcal X$, which is
problem-specific. Note the reduction is \emph{black-box}: line~7 may be replaced
by any $\gamma$-approximate minimizer of $F(\Chat,\cdot)$, yielding overall ratio
$\gamma\cdot\frac{1+\varepsilon}{1-\varepsilon}$ (Theorem~\ref{thm:approxalg}).
\end{document}
